% Section 4 — Information parity and the irrelevance theorem
% (charter §4.4; issues T1-2, T2-1, T2-2)
\section{Information parity and the irrelevance theorem}
\label{sec:parity}

% STUB (T1-2, T2-1, T2-2). The spine of the paper.
% Area role (charter §2): parity = sufficiency / Markov-equivalence; the
% irrelevance theorem via DPI; rate-distortion for lossy reps.
%
% Plan:
%   - T2-1: formal definitions of structured representation and information
%     parity (parity = \rep(X) \suff X w.r.t. Y, i.e. X \markov \rep(X) \markov Y
%     with no loss of sufficiency).
%   - T2-2: the at-optimum invariance theorem: under parity, \bayesrisk is
%     invariant to \rep. DPI is the engine.
%   - Corollary: rate-distortion bound \ratedist for lossy representations.

\noindent\emph{[Stub. Definitions in T2-1; irrelevance theorem in T2-2;
information-theoretic foundations in T1-2.]}

\begin{theorem}[Irrelevance at optimum --- placeholder]
\label{thm:irrelevance}
\emph{[Placeholder for T2-2.]} Under information parity
($X \markov \rep(X) \markov Y$ with $\rep(X) \suff X$), the Bayes risk is
format-invariant: $\bayesrisk(\rep) = \bayesrisk(\mathrm{id})$.
\idealization{This statement assumes Bayes-optimality (the
infinite-capacity limit) and exact parity. A real frozen LLM satisfies
neither; \cref{sec:parity-defect} quantifies the parity defect and
\cref{sec:bounded} relaxes the capacity idealization. ``Parity is
operational, not absolute'' (preprint §3.3) is carried in here as a formal
caveat.}
\end{theorem}
